\input eplain
\pdfpagewidth=148mm
\pdfpageheight=210mm
\hsize=128mm
\vsize=190mm

\pdfhorigin=0pt
\pdfvorigin=0pt
\hoffset=10mm
\voffset=10mm
\noindent news \leaders\hbox to 10pt {\hss.\hss}  \hfill \xref{node2}
\smallskip
\noindent 02 limits \leaders\hbox to 10pt {\hss.\hss}  \hfill \xref{node0}
\smallskip
\noindent 04 infinite limits \leaders\hbox to 10pt {\hss.\hss}  \hfill \xref{node1}
\smallskip
\vfill \break
{\medskip\noindent\bf\font\titlefont=cmss17 at 20pt\titlefont Timeline}\medskip
{\noindent \bf 2026-01-25:} 02 limits\smallskip
{\noindent \bf 2026-01-26:} 02 limits\smallskip
{\noindent \bf 2026-01-27:} 02 limits\smallskip
{\noindent \bf 2026-01-28:} news, 02 limits, 04 infinite limits\smallskip
{\noindent \bf 2026-01-29:} 02 limits\smallskip
\vfill \break
\xrdef{node0}
{\medskip\noindent\bf\font\titlefont=cmss17 at 20pt\titlefont 02 limits}\medskip
{\medskip\noindent\bf\font\titlefont=cmitt12 at 17pt\titlefont 2026-01-25}\medskip
{\smallskip\noindent\bf\font\titlefont=cmr12 at 13pt\titlefont $\bullet$ 13:48: Calculus chapter 2: 2.1, 2.2, 2.3}\smallskip
Chapter 2: limits
\par

2.1: the idea of limits. The whole point of this chapter was to tie the relationship of average velocity and instantanous velocity to the relationship of the slope of a secant line to a tangent.
\par

2.2: definitions of limits. a preliminary definition of a limit is provided. one-sided limits (left/right). a theorem is provided that states that both left/right limits must exist for a 2-sided limit to exist.
\par

2.3 techniques for computing limits. limits of linear functions. limit laws.
\par

{\smallskip\noindent\bf\font\titlefont=cmr12 at 13pt\titlefont $\bullet$ 15:45: some more 2.3 probably}\smallskip
{\smallskip\noindent\bf\font\titlefont=cmr12 at 13pt\titlefont $\bullet$ 16:17: wrap up 2.3}\smallskip
proof of power law using product law. finding limits of polynomials and rational functions, with proofs using limit laws.
\par

{\medskip\noindent\bf\font\titlefont=cmitt12 at 17pt\titlefont 2026-01-26}\medskip
{\smallskip\noindent\bf\font\titlefont=cmr12 at 13pt\titlefont $\bullet$ 11:35: more 2.3 actually}\smallskip
Looked at examples for limits of rational functions. one sided limits, and an adjusted root law. Other techniques for finding limits, which involved finding equivalent limits where direct subsitution could be used.
\par

{\smallskip\noindent\bf\font\titlefont=cmr12 at 13pt\titlefont $\bullet$ 13:15: 2.3}\smallskip
{\smallskip\noindent\bf\font\titlefont=cmr12 at 13pt\titlefont $\bullet$ 15:26: 2.3}\smallskip
{\smallskip\noindent\bf\font\titlefont=cmr12 at 13pt\titlefont $\bullet$ 16:03: 2.3}\smallskip
{\medskip\noindent\bf\font\titlefont=cmitt12 at 17pt\titlefont 2026-01-27}\medskip
{\smallskip\noindent\bf\font\titlefont=cmr12 at 13pt\titlefont $\bullet$ 16:35: review: complex conjugates, factoring (quickly)}\smallskip
One of the examples in the textbook uses a conjugate to find and equivalent limit that to apply limit laws with. I completely forgot what a conjugate was, but found a wiki page on the subject (specifically related to square roots).
\par

It turns out we are factoring, so I had to run home and remember some of the tricks involved in factoring quadratics. It's coming back.
\par

{\smallskip\noindent\bf\font\titlefont=cmr12 at 13pt\titlefont $\bullet$ 18:28: more 2.3 (techniques for computing limits)}\smallskip
Limits that do not work with limit law techniques. Exponentials were used as an example. squeeze theorem. also some work with trig limits, for which squeeze theorem was used.
\par

{\smallskip\noindent\bf\font\titlefont=cmr12 at 13pt\titlefont $\bullet$ 19:50: more 2.3}\smallskip
{\medskip\noindent\bf\font\titlefont=cmitt12 at 17pt\titlefont 2026-01-28}\medskip
{\smallskip\noindent\bf\font\titlefont=cmr12 at 13pt\titlefont $\bullet$ 10:20: some problem sets from chapter 2}\smallskip
Had some issues accessing the HW today, so I took some time this morning to do some problems from the book.
\par

I've also begun a "scratch" index system. Basically, I'm going to hold onto and index the scraps I work out problems on. Might be useful dunno yet.
\par

{\medskip\noindent\bf\font\titlefont=cmitt12 at 17pt\titlefont 2026-01-29}\medskip
{\smallskip\noindent\bf\font\titlefont=cmr12 at 13pt\titlefont $\bullet$ 09:48: Calculus reading: initial reading on 2.5 (limits at infinity)}\smallskip
\xrdef{node1}
{\medskip\noindent\bf\font\titlefont=cmss17 at 20pt\titlefont 04 infinite limits}\medskip
{\medskip\noindent\bf\font\titlefont=cmitt12 at 17pt\titlefont 2026-01-28}\medskip
{\smallskip\noindent\bf\font\titlefont=cmr12 at 13pt\titlefont $\bullet$ 16:45: 2.4 infinite limits}\smallskip
Infinite limits vs limits at infinity. finding limits analytically. example 4 was tricky to work out for me, especially with my factoring being rusty. Slowly figuring out when things go to negative infinity or positive infinity for one-sided limits based on approach.
\par

{\smallskip\noindent\bf\font\titlefont=cmr12 at 13pt\titlefont $\bullet$ 18:12: 2.4 infinite limits cont.}\smallskip
\xrdef{node2}
{\medskip\noindent\bf\font\titlefont=cmss17 at 20pt\titlefont news}\medskip
{\medskip\noindent\bf\font\titlefont=cmitt12 at 17pt\titlefont 2026-01-28}\medskip
{\smallskip\noindent\bf\font\titlefont=cmr12 at 13pt\titlefont $\bullet$ 18:59: logging, seting up some new things related to my calculus class}\smallskip
\bye
